\documentclass[12pt,letterpaper]{article}
\usepackage[margin=1in]{geometry}
\usepackage[T1]{fontenc}
\usepackage[utf8]{inputenc}
\usepackage[french]{babel}
\usepackage{graphicx}
\usepackage{url}
\usepackage{etoolbox}
\usepackage{xcolor}
\usepackage{amsmath}
\usepackage{amssymb}
\usepackage{float}
\usepackage{listings}
\usepackage{listingsutf8}
\usepackage{verbatim}

% Code listing style for Python
\lstdefinestyle{pythonstyle}{
  language=Python,
  inputencoding=utf8/latin1,
  basicstyle=\ttfamily\small,
  keywordstyle=\color{blue!70!black}\bfseries,
  commentstyle=\color{green!40!black}\itshape,
  stringstyle=\color{orange!70!black},
  numbers=left,
  numberstyle=\scriptsize\color{gray!70},
  stepnumber=1,
  numbersep=10pt,
  frame=single,
  framerule=0.3pt,
  rulecolor=\color{gray!40},
  backgroundcolor=\color{gray!5},
  breaklines=true,
  breakatwhitespace=true,
  showstringspaces=false,
  tabsize=4,
  columns=fullflexible,
  keepspaces=true,
  xleftmargin=12pt,
  xrightmargin=6pt,
  aboveskip=8pt,
  belowskip=8pt,
  captionpos=b
}

\lstset{style=pythonstyle}

% Start each section on a new page
\pretocmd{\section}{\clearpage}{}{}

\title{MTH8302~-- Devoir 3: R\'{e}gression Lin\'{e}aire Multiple}
\author{Leonard Excoffier (2085276)}
\date{Hiver 2026}

\begin{document}

\maketitle

\tableofcontents

\section*{Introduction}

Le d\'{e}p\^{o}t GitHub~\cite{repo} contient le code complet, structur\'{e} et pr\^{e}t \`{a} \^{e}tre ex\'{e}cut\'{e}.
Tous les scripts Python se trouvent dans \texttt{src/} et les sorties (figures, fichiers texte) dans \texttt{output/}.
Le script \texttt{run.sh} ex\'{e}cute l'ensemble des scripts puis compile le rapport automatiquement.
Le PDF compil\'{e} est g\'{e}n\'{e}r\'{e} \`{a} la racine du projet.

\begin{thebibliography}{9}
\bibitem{repo}
Leonard Excoffier.~\emph{MTH8302~-- Devoir 3: R\'{e}gression Lin\'{e}aire Multiple}.\ GitHub repository.
\url{https://github.com/excoffierleonard/mth8302-devoir3.git}. Consult\'{e} le 9 mars 2026.
\end{thebibliography}

\end{document}
